\section{Discusi\'on}

En relaci\'on con el primer objetivo, la discretizaci\'on conserv\'o la potencia integrada de ambas fuentes y produjo una matriz sim\'etrica definida positiva, compatible con los dos solucionadores. Para el segundo objetivo, la Figura~\ref{fig:resultados-base} muestra el m\'aximo dentro de la huella GPU, cuya potencia base fue mayor que la de CPU. La Tabla~\ref{tab:base-solver} cuantifica la respuesta t\'ermica, mientras que estudios de placas reales incorporan capas, interacciones entre componentes y mecanismos de disipaci\'on que aqu\'i se excluyeron \citep{harvest2007devices,moniervinard2015pcb,zhang2024pcb,skadron2003hotspot}. En consecuencia, $77.759239\ ^\circ$C es un resultado del modelo idealizado, no una temperatura operativa medida.

El cuarto objetivo se respondi\'o comparando ambos m\'etodos con la misma $A$, $\mathbf{b}$, condici\'on inicial y tolerancia. La Tabla~\ref{tab:base-solver} y la curva de residuos de la Figura~\ref{fig:resultados-base} muestran que CG alcanz\'o $10^{-8}$ en 127 iteraciones, frente a 2145 de Gauss--Seidel. En la malla m\'as fina, los conteos fueron 485 y 29\,541, respectivamente. La Figura~\ref{fig:estudios-completos} presenta medianas e intervalos intercuart\'ilicos medidos tras calentamiento y en tres repeticiones por configuraci\'on; en esas ejecuciones CG requiri\'o menos tiempo, aunque la dispersi\'on y la dependencia del equipo limitan la generalizaci\'on de los cronometrajes. Esta comparaci\'on concuerda con el uso de CG en sistemas sim\'etricos definidos positivos y con antecedentes de m\'etodos iterativos para Poisson, pero los resultados propios dependen de la matriz, el ordenamiento y la implementaci\'on utilizados \citep{zhang2001,xu2011poisson,oladejo2007tool}.

La Figura~\ref{fig:aproximaciones} muestra que los campos iterativos se construyeron desde el estado inicial uniforme de $25\ ^\circ$C hasta las soluciones convergidas; los paneles intermedios corresponden a iteraciones algebraicas y no a tiempos de calentamiento. La Figura~\ref{fig:resultados-base} permite localizar las fuentes, inspeccionar la geometr\'ia del campo y comparar la reducci\'on del residuo; sus paneles representan el mismo sistema y no constituyen evidencia independiente.

Respecto del tercer objetivo, la Tabla~\ref{tab:refinamiento} muestra que al duplicar la resoluci\'on las diferencias m\'aximas entre soluciones de mallas consecutivas disminuyeron de $0.23847$ a $0.06159$ y $0.01554\ ^\circ$C. La tendencia, cercana a una reducci\'on por cuatro, es compatible con el comportamiento esperado de diferencias centrales de segundo orden en las mallas evaluadas; no equivale a un error exacto porque no se dispone de soluci\'on anal\'itica para las fuentes discontinuas y la malla m\'as fina carece de una malla de referencia superior. La Figura~\ref{fig:estudios-completos} completa el an\'alisis al mostrar el aumento del costo computacional al crecer el sistema. La interpretaci\'on es consistente con an\'alisis de esquemas para Poisson y trabajos de refinamiento, sin atribuir sus garant\'ias te\'oricas a esta prueba particular \citep{settle2013compact,chen2007supra,raeli2018hierarchical,batty2017cell}.

La Figura~\ref{fig:estudios-completos} resume las ocho configuraciones de fuentes y muestra un aumento del m\'aximo al elevar por separado la potencia de CPU o GPU. La ecuaci\'on discreta es lineal respecto de las potencias cuando se mantienen las propiedades y las fronteras; sin embargo, $T_{\max}$ es el m\'aximo global del campo y puede cambiar de ubicaci\'on, por lo que la curva no tiene que ser una recta con pendiente constante. La sensibilidad orienta el an\'alisis de esta idealizaci\'on y no sustituye los modelos multicapa, transitorios y con disipaci\'on que se aplican a placas reales \citep{chung2001materials,zheng2021advances,harvest2007devices,moniervinard2015pcb}.

La verificaci\'on complementaria confirm\'o residuo relativo menor que la tolerancia, concordancia cercana con la referencia directa, simetr\'ia, definitud positiva, cumplimiento de frontera y conservaci\'on de la potencia discretizada. Las pruebas de tolerancia, potencia extrema y potencia negativa evaluaron el comportamiento algebraico; las dos pruebas de potencia no describen condiciones f\'isicas. La diferencia entre soluci\'on iterativa y referencia directa estima el error algebraico, mientras que la comparaci\'on entre mallas caracteriza el cambio espacial. Ninguna mide el error de modelaci\'on f\'isica, que requerir\'ia comparar con una placa y condiciones experimentales documentadas. Las soluciones de verificaci\'on para s\'olidos con calentamiento parcial y los modelos num\'ericos alternativos de conducci\'on proveen referentes metodol\'ogicos, pero sus geometr\'ias y condiciones de frontera no validan directamente THERMOPLATE \citep{beck2004verification,akeremale2020finite,reutskiy2016meshless,deng2017high}.
